\section{Introducción}
La Red Dorsal Nacional de Fibra Óptica (RDNFO) es un proyecto que tiene como finalidad el despliegue y
funcionamiento de una red de fibra óptica que conectará la capital del Perú con diferentes capitales de región y capitales de 
provincia, en este mismo sentido el diseño permitirá la reducción de los costos para que los habitantes de zonas apartadas puedan 
acceder a una conexión de banda ancha, en este mismo sentido 
permitirá que ellos puedan contar con una conexión estable a internet.\\
Sin embargo, el principal problema de este proyecto es que, a pesar de estar instalada casi en su totalidad, el porcentaje de 
operación actual es mínimo. Esta situacion ha generado un estancamiento que ha afectado críticamente al país durante la pandemia, 
permitiendo que las redes de telecomunicaciones privadas 
capturen la mayor parte del mercado y retrasando la inclusión digital en varias comunidades y zonas aledañas. Por lo tanto, el análisis de los 
fallos de la RDNFO son clave para promover el desarrollo social del Perú, esto porque los gastos de mantenimiento de semejante infraestructura 
subutilizada representa una perdida de recursos públicos.\\
Ante este escenario, el objetivo general de la presente monografía es analizar las causas técnicas, económicas y contractuales que han llevado al 
estancamiento y subutilización de la Red Dorsal Nacional de Fibra Óptica, evaluando su impacto en la conectividad del país.\\